%% ---------------------------------------------------------------------------
%% Manuscript v4 (edited from v3; audit-study scope) -- status header
%%   NOT ICML-READY. Learned-binder method development is on HOLD.
%%   TARGET_YEAR        = 2027   (ICML 2027 style NOT published: icml.cc 2027
%%                                 pages returned 404 on 2026-09-26)
%%   TEMPLATE_YEAR      = 2026   (temporary; official ICML 2026 style)
%%   SUBMISSION_READY   = false
%%   STYLE_FILES_PRESENT= false  (icml2026.sty / icml2026.bst are NOT in this
%%                                 repository; download of icml2026.zip was not
%%                                 approved. Do not replace them with a
%%                                 hand-made imitation.)
%%   COMPILE_STATUS     = COMPILE_NOT_RUN (no TeX installation in the
%%                        environment; see BUILD.md for exact commands)
%%   REVIEW MODE        = \usepackage{icml2026} without [accepted]. This is
%%                        inferred from the ICML 2026 author instructions,
%%                        which name [accepted] for camera-ready; the example
%%                        paper in icml2026.zip was not read. Check the title
%%                        block macros below against it before compiling.
%%   PAGE LIMIT         = 8 pages main body at submission (ICML 2026 author
%%                        instructions); impact statement, references and
%%                        appendices do not count. Length NOT verified by a
%%                        compile; v3 moves the simulator study and V1 census
%%                        details to the appendix (word count in STATUS.md).
%%   SCOPE              = oracle-plan audit (title and abstract); strict and
%%                        fallback arms reported separately; no language
%%                        model, no learning.
%%   All numbers come from paper/tables/*.tex, generated by
%%   scripts/make_paper_tables.py from the files listed in paper/sources.json.
%% ---------------------------------------------------------------------------
\documentclass{article}

\usepackage{microtype}
\usepackage{graphicx}
\usepackage{booktabs}
\usepackage{hyperref}
\usepackage{xcolor}

\usepackage{icml2026}

\usepackage{amsmath}
\usepackage{amssymb}
\usepackage{amsthm}

\newtheorem{definition}{Definition}
\newtheorem{proposition}{Proposition}
\newtheorem{observation}{Observation}
\newcommand{\todo}[1]{{\color{red}[TODO: #1]}}
\newcommand{\notrun}{\textsc{not run}}
\newcommand{\abst}{\bot}

% generated by scripts/make_paper_tables.py from paper/sources.json -- do not edit
\newcommand{\EOneRows}{15{,}750}
\newcommand{\EOneWall}{5.4}
\newcommand{\EOneCommit}{54caffb}
\newcommand{\RuleSS}{0.773}
\newcommand{\LearnSS}{0.798}
\newcommand{\CamelSS}{0.611}
\newcommand{\DocSS}{0.560}
\newcommand{\DocRefViol}{0.067}
\newcommand{\RuleASR}{0.178}
\newcommand{\LearnASR}{0.194}
\newcommand{\DiffTemplate}{+0.024}
\newcommand{\DiffEnv}{-0.007}
\newcommand{\DiffInstance}{-0.033}
\newcommand{\DiffInstanceAtk}{+0.032}
\newcommand{\LexAgreeTotal}{1920}
\newcommand{\ArgsTotal}{1920}
\newcommand{\LearnAgreeTotal}{1505}
\newcommand{\HSim}{0.199}
\newcommand{\ProxySim}{1.000}
\newcommand{\GuardedRefViol}{0}
\newcommand{\GuardedLeak}{0}
\newcommand{\LearnOutOfScopeProposals}{28}
\newcommand{\AgentDojoStatus}{NOT\_RUN}
\newcommand{\AdjUser}{97}
\newcommand{\AdjResolved}{89}
\newcommand{\AdjUnresolved}{8}
\newcommand{\AdjZeroSlot}{52}
\newcommand{\AdjVariants}{940}
\newcommand{\AdjInjVariants}{851}
\newcommand{\AdjHlo}{0.253}
\newcommand{\AdjHup}{0.335}
\newcommand{\AdjHfirstLo}{0.216}
\newcommand{\AdjHfirstUp}{0.299}
\newcommand{\AdjHsecLo}{0.306}
\newcommand{\AdjHsecUp}{0.306}
\newcommand{\AdjBenignHlo}{0.247}
\newcommand{\AdjInjHlo}{0.254}
\newcommand{\AdjRuleAtk}{4}
\newcommand{\AdjGtAtk}{0}
\newcommand{\AdjRemoved}{34}
\newcommand{\AdjAtkAdmitted}{37}
\newcommand{\AdjWall}{345.7}
\newcommand{\AdjNonExh}{0}
\newcommand{\AdjUnsolvable}{9}
\newcommand{\AdjUnsolvableTasks}{1}
\newcommand{\AdjNumKept}{81}
\newcommand{\AdjNumLo}{0.179}
\newcommand{\AdjNumUp}{0.228}
\newcommand{\AdjNumericSlots}{18}
\newcommand{\AdjSlots}{91}
\newcommand{\CcHLo}{0.273}
\newcommand{\CcHUp}{0.293}
\newcommand{\CcHRuleLo}{0.314}
\newcommand{\CcHRuleUp}{0.334}
\newcommand{\CcHEffLo}{0.273}
\newcommand{\CcHEffUp}{0.293}
\newcommand{\CcHTaskLo}{0.273}
\newcommand{\CcHTaskUp}{0.293}
\newcommand{\CcHRuleTaskLo}{0.314}
\newcommand{\CcHRuleTaskUp}{0.334}
\newcommand{\CcHStrictLo}{0.271}
\newcommand{\CcHStrictUp}{0.292}
\newcommand{\CcHRuleStrictLo}{0.302}
\newcommand{\CcHRuleStrictUp}{0.323}
\newcommand{\CcHBenLo}{0.289}
\newcommand{\CcHBenUp}{0.309}
\newcommand{\CcHInjLo}{0.270}
\newcommand{\CcHInjUp}{0.291}
\newcommand{\CcTasks}{97}
\newcommand{\CcResolved}{97}
\newcommand{\CcUnresolved}{0}
\newcommand{\CcStrictResolved}{97}
\newcommand{\CcVariants}{1{,}046}
\newcommand{\CcInjVariants}{949}
\newcommand{\CcNoWrite}{36}
\newcommand{\CcIncomplete}{77}
\newcommand{\CcLeaves}{340{,}667}
\newcommand{\CcResAtk}{4}
\newcommand{\CcRuleAtk}{4}
\newcommand{\CcAbsAtk}{0}
\newcommand{\CcGtAtk}{0}
\newcommand{\CcGtN}{709}
\newcommand{\CcResAtkEff}{4}
\newcommand{\CcRuleAtkEff}{4}
\newcommand{\CcReachLo}{35}
\newcommand{\CcReachUp}{110}
\newcommand{\CcReachEffLo}{35}
\newcommand{\CcReachEffUp}{110}
\newcommand{\CcReachTnLo}{0.044}
\newcommand{\CcReachTnUp}{0.119}
\newcommand{\CcSelN}{248}
\newcommand{\CcSelTasks}{28}
\newcommand{\CcSelUtil}{248}
\newcommand{\CcSelAtk}{0}
\newcommand{\CcSelAbst}{165}
\newcommand{\CcNotRefuted}{248}
\newcommand{\CcRefuted}{0}
\newcommand{\CcIdUnknown}{0}
\newcommand{\CcSingletonView}{148}
\newcommand{\CcFirstWriteErr}{130}
\newcommand{\CcLaterErr}{70}
\newcommand{\CcNRRate}{0.273}
\newcommand{\CcOutDep}{42}
\newcommand{\CcOracleComposed}{6}
\newcommand{\CcFreeText}{46}
\newcommand{\CcGoldNotAdmitted}{21}
\newcommand{\CcGtFails}{17}
\newcommand{\CcPlanVaries}{0}
\newcommand{\CcWall}{2183}
\newcommand{\CcCPU}{4100}
\newcommand{\CcFollowup}{PROPOSE\_CONDITIONAL\_NOT\_RUN}
\newcommand{\CcSelViewShared}{100}
\newcommand{\CcUnkLocus}{48}
\newcommand{\CcAbstFirstWrite}{102}
\newcommand{\CcGapTasks}{28}
\newcommand{\CcUndetTasks}{3}
\newcommand{\CcWriteTasks}{61}
\newcommand{\CcBenUnsolv}{14}
\newcommand{\CcChanDiv}{0}
\newcommand{\CcPathRecs}{4{,}184}
\newcommand{\CcToolErrPaths}{31}
\newcommand{\CcOutDepWrite}{9}
\newcommand{\CcOutDepCompOverlap}{4}
\newcommand{\CcOutDepCompUnion}{44}
\newcommand{\CcSOutDep}{38}
\newcommand{\CcSecUndet}{30}
\newcommand{\CcSecUndetTasks}{3}
\newcommand{\CcIncompleteTasks}{9}
\newcommand{\CcAtkUndet}{75}
\newcommand{\CcSResAtk}{0}
\newcommand{\CcSRuleAtk}{0}
\newcommand{\CcSReachLo}{31}
\newcommand{\CcSReachUp}{106}
\newcommand{\CcSGtN}{674}
\newcommand{\CcSSelN}{247}
\newcommand{\CcSSelTasks}{27}
\newcommand{\CcVacuousTasks}{15}
\newcommand{\CcCovTasks}{21}
\newcommand{\CcCovSlots}{43}
\newcommand{\CcCovNotInView}{29}
\newcommand{\CcCovInText}{12}
\newcommand{\CcCovList}{2}
\newcommand{\CcCovDatetime}{18}


\icmltitlerunning{Evidence Binding under Fixed Authority Boundaries: An Oracle-Plan Audit}

\begin{document}

\twocolumn[
\icmltitle{Evidence Binding under Fixed Authority Boundaries\\in Tool-Using Agents: An Oracle-Plan Audit on AgentDojo}

\begin{icmlauthorlist}
\icmlauthor{Anonymous Authors}{anon}
\end{icmlauthorlist}
\icmlaffiliation{anon}{Anonymous Institution}
\icmlcorrespondingauthor{Anonymous Authors}{anonymous@example.com}
\icmlkeywords{prompt injection, tool-using agents, provenance, evidence binding}

\vskip 0.3in
]

\printAffiliationsAndNotice{}

%% ===========================================================================
\begin{abstract}
This is an audit under an \emph{oracle plan}. On AgentDojo v1.2.2 we hold
the benchmark's ground-truth call sequence, read arguments, free-text
arguments and final answer fixed. A binder may only pick one typed
candidate admitted by a fixed contract, or abstain. Every admitted choice
sequence is replayed offline against the benchmark's own evaluators, without
a language model.

Slot types, candidates and the baselines' decisions are derived from the
tool schema, the request and earlier observations; they do not change when
hidden gold values are mutated. The evaluated arms are still
oracle-conditioned. The \emph{fallback} arm additionally fills slots that
have no observed candidate with the planner's gold value; the \emph{strict}
arm abstains there instead.

Over a typed lexical resolver, the clairvoyant envelope of what any admitted
choice could add is $[\CcHStrictLo,\CcHStrictUp]$ (strict) and
$[\CcHLo,\CcHUp]$ (fallback). This is a computational envelope over
partially searched trees, not a confidence interval or an achievable gain.
In both arms the utility-only envelope is identical: every selectable
failure is a utility failure, e.g.\ choosing the channel with the most
users.

The baselines executed the attacker goal in \CcSResAtk{} (strict) and
\CcResAtk{} (fallback) of \CcInjVariants{} injected variants. Some admitted
choice achieves it in \CcSReachLo--\CcSReachUp{} and \CcReachLo--\CcReachUp{}
variants respectively. These are witness counts, not attack rates, and they
show that risk within the contract is not zero.

We do not pursue a learned binder on this evidence. The audit does not show
that argument binding is unimportant for agents with real planners.
\end{abstract}



%% ===========================================================================
\section{Introduction}
\label{sec:intro}

Tool-using language-model agents read e-mails, files and web pages written by
third parties, then act on the user's behalf. Indirect prompt injection turns
that content into an attack surface. A line of defenses therefore separates
what the agent is allowed to do from what it reads.

\begin{itemize}
\item \textbf{Control/data separation.} CaMeL derives a program from the
trusted query and tags every value with capabilities that policies check
before each tool call \citep{debenedetti2025camel}. Information-flow systems
propagate confidentiality and integrity labels
\citep{wu2024fsecure,costa2025fides,zhong2025rtbas}. Privilege-control
systems restrict tools and arguments \citep{shi2025progent,tsai2025conseca}.
\item \textbf{Argument-level provenance.} The most recent systems move to the
granularity of individual arguments. PACT checks each argument's provenance
against a role-specific contract \citep{fan2026pact}. ROPE admits a
sensitive parameter only if its origin satisfies a marker chosen from the
user's request \citep{ma2026rope}. ARGUS releases an action only when benign
evidence entails every argument \citep{weng2026argus}.
\end{itemize}

\paragraph{Authorization is not correctness.}
These boundaries decide whether a value is \emph{admissible}. ROPE states the
limit plainly: its guarantee ``is about origin, not intent: it ensures an
admitted value's origin is trusted, not that it is the value the user
wanted'' (\citealp{ma2026rope}, \S3.4). Consider ``pay the invoice for order
ORD-1234 from our vendor''. Once the vendor is delegated as a source, all of
the following are admissible:
\begin{itemize}
\item the current invoice;
\item an older invoice for a different order;
\item a draft that was later corrected;
\item a message from the vendor's compromised account announcing ``new bank
details''.
\end{itemize}
Some component must still choose among admissible values, or abstain. In
current systems that component is the planner language model, a hand-written
rule, or a prompted grounding step. ARGUS notes that ``a forged invoice can
make the attacker's account the only available evidence''
(\citealp{weng2026argus}, \S6.2).

\paragraph{Question and scope.}
We ask how much this choice can matter when everything else is held fixed:
the plan, the per-argument authority contract, and the information available
when the argument is bound. We answer only under an oracle plan on
AgentDojo, and we train nothing.

\paragraph{What this paper reports.}
Three census findings, each conditional on the oracle plan
(\S\ref{sec:v2-results}):
\begin{enumerate}
\item \textbf{Utility, not security.} The envelope over a typed resolver is
$[\CcHStrictLo,\CcHStrictUp]$ (strict) and $[\CcHLo,\CcHUp]$ (fallback), and
it is identical for utility alone. Every selectable failure is a choice
among legitimate values.
\item \textbf{Rare execution, non-zero admitted risk.} Baselines rarely
executed the attacker goal, but some admitted choice achieves it in a
bounded set of injected variants.
\item \textbf{Limits of the measurement.} Some tasks cannot be completed by
any candidate our generator admits. The identifiability screen is vacuous
for most selectable failures.
\end{enumerate}
The formulation (\S\ref{sec:setup}) and the ceiling
(\S\ref{sec:headroom}) are measurement framing built on prior contract and
provenance work; we do not count them as contributions. Our earlier census
(retracted as a deployable ceiling, \S\ref{sec:v1}) and a synthetic
simulator study (Appendix~\ref{sec:sim}) are kept as a record. They are not
contributions either.

\paragraph{Non-claims.}
We do not claim the following:
\begin{itemize}
\item that learned binding improves security or utility of real agents;
\item that the AgentDojo ceiling transfers to agents whose planner is a
language model rather than the ground truth;
\item that the V1 census measures a deployable contract;
\item that the selectable failures are identifiable from observation (they
are only not refuted);
\item that the evaluated arms are gold-free (only the actor-view derivation
is), or that the official and executed-effect channels are equivalent in
general;
\item that the admitted-attack witness counts are attack rates or
learnable security gains;
\item that tasks our candidate generator cannot serve are out of reach for
every binder;
\item that our validator provides process isolation or a formal guarantee;
\item that any simulator number measures a real model;
\item that ``span binding'' alone distinguishes this work from ROPE, PACT or
ARGUS (\S\ref{sec:related}).
\end{itemize}

%% ===========================================================================
\section{Related Work}
\label{sec:related}

\paragraph{System-level defenses.}
The dual-LLM and plan-then-execute patterns keep untrusted data away from the
component that chooses actions \citep{beurerkellner2025design}.
\begin{itemize}
\item \textbf{CaMeL} \citep{debenedetti2025camel}: a privileged model writes
a program from the trusted query, and a quarantined model parses untrusted
data. An interpreter tags each value with its data provenance and allowed
readers and evaluates Python security policies before tool calls; tools can
also report an ``inner source'', such as an e-mail sender.
\item \textbf{f-secure} \citep{wu2024fsecure} and \textbf{FIDES}
\citep{costa2025fides}: information-flow control for agents.
\item \textbf{RTBAS} \citep{zhong2025rtbas}: adds a dependency screener,
including a trained one, that selects which context regions an action
depends on.
\item \textbf{Permissive IFC} \citep{siddiqui2024permissive}: propagates only
the labels of influential inputs.
\item \textbf{AirGapAgent} \citep{bagdasarian2024airgap}: minimizes the data
exposed for a task.
\item \textbf{Progent} \citep{shi2025progent} and \textbf{Conseca}
\citep{tsai2025conseca}: express privilege policies that a model may help
write and that are enforced deterministically.
\item \textbf{Jacob et al.} \citep{jacob2025typed}: restrict what crosses the
privilege boundary to narrow data types.
\end{itemize}
Detection- and alignment-style defenses include MELON
\citep{zhu2025melon}, Task Shield \citep{jia2024taskshield} and IPIGuard
\citep{an2025ipiguard}. Our study takes all of these mechanisms as given. We
claim none of control/data separation, provenance or capability enforcement,
typed restriction, or learned context selection as new.

\paragraph{Argument-level authority.}
Three recent systems are closest to our setting (Table~\ref{tab:rw}).
\begin{itemize}
\item \textbf{PACT} \citep{fan2026pact} assigns each argument a semantic role
and a contract $C_t$. It tags values with provenance
$\pi(v)=\langle\mathcal{O}(v),\tau(v),\mathcal{B}(v)\rangle$ and admits an
argument only if its trust level, origins and obligations satisfy the
contract. It reports that, with oracle provenance, contract checking reaches
full utility and security on its diagnostic suites. It locates ``the
remaining deployment bottleneck'' in provenance inference and contract
synthesis.
\item \textbf{ROPE} \citep{ma2026rope} checks sensitive parameters against
origin markers (e.g.\ \emph{prompt}, \emph{sourced}, \emph{record}). A
zero-shot language-model router reads only the user's request and outputs
the scope. An ``enforcement clamp'' accepts a router's override only when it
is at least as strict as an offline default. Origins are assigned from the
platform's own division of a result, so ``a file body or a web page \ldots\
arrives as one block'' (\S3.3).
\item \textbf{ARGUS} \citep{weng2026argus} builds an influence-provenance
graph. It segments context into spans labelled benign or anomalous, and
grounds each argument in supporting spans with prompted sub-agents.
\end{itemize}

\paragraph{How we differ.}
The difference is narrow. Compared with each system:
\begin{itemize}
\item \textbf{ROPE}'s router selects a \emph{scope}; it does not choose among
several admitted values.
\item \textbf{PACT}'s check decides admissibility, not which admissible value
is correct.
\item \textbf{ARGUS} does ground arguments in spans, but with prompted models
and a release decision rather than an abstaining binder under an external
fixed contract.
\end{itemize}
What we add is a way to \emph{measure} how much the choice among admitted
values can matter while the contract is fixed. Span-level grounding itself
is not new (ARGUS), and we do not rest novelty on it. Whether a \emph{trained}
binder beats both rules and prompted grounding at matched cost is an open
empirical question that this paper does not settle.

\paragraph{Benchmarks.}
AgentDojo \citep{debenedetti2024agentdojo} provides user tasks, injection
tasks and post-state utility and security checks. PACT evaluates on
``all 97 benign user tasks and all 27 injection tasks''
\citep{fan2026pact}. AgentDyn \citep{li2026agentdyn} adds ``60 challenging
open-ended tasks and 560 injection test cases'' that include helpful
third-party instructions. We replay AgentDojo offline without a language
model (\S\ref{sec:exp-real}); we have not run AgentDyn.

\begin{table*}[t]
\caption{Closest systems, from the sections we read (Appendix~\ref{app:reading}).
``Chooser'' is the component that picks among several \emph{admitted}
values. None of the listed systems trains its chooser; RTBAS trains a
dependency screener, not an argument binder.}
\label{tab:rw}
\centering\small
\begin{tabular}{lllll}
\toprule
System & Authority unit & Provenance granularity & Chooser among admitted values & Evaluation \\
\midrule
CaMeL & per value, tool-call policies & value; tool ``inner source'' & program + quarantined LLM (prompted) & AgentDojo \\
PACT & per argument role contract & value tag $\langle\mathcal{O},\tau,\mathcal{B}\rangle$ & agent LLM (not addressed) & AgentDojo \\
ROPE & per sensitive parameter & platform-divided entry & agent LLM; router picks scope only & AgentDojo, AgentDyn \\
ARGUS & per action (release/block) & LLM-segmented spans & prompted argument grounder & AgentLure, AgentDojo \\
This paper & per typed argument, fixed contract & actor-view value & rules, typed resolver; ceiling only & AgentDojo replay; sim.\ (app.) \\
\bottomrule
\end{tabular}
\end{table*}

%% ===========================================================================
\section{Problem Setup}
\label{sec:setup}

\subsection{Authorization versus factual correctness}
A task $t$ consists of a trusted user instruction $q_t$, a \emph{fixed plan}
of tool calls, and a set of argument slots $K_t$ of state-changing calls.
For each variant $v$ of the task (benign, or one injection), the runtime
reads tool outputs and extracts a candidate set $X_{t,v}$.
\begin{itemize}
\item Each candidate $x$ carries a value $\mathrm{val}(x)$ and a provenance
record $\pi(x)$.
\item A \emph{contract} $\kappa_t$ assigns each slot $k$ a predicate
$\kappa_{t,k}$ on provenance records.
\item The \emph{admitted set} is
$C_{t,v,k}=\{x\in X_{t,v}:\kappa_{t,k}(\pi(x))\}$.
\end{itemize}
An assignment $a$ gives each slot an element of $C_{t,v,k}\cup\{\abst\}$,
where $\abst$ means abstain. A call is \emph{authorized} if every slot is
bound to an admitted candidate. The task \emph{succeeds},
$s(t,v,a)=1$, if the post-state after executing the fixed plan with $a$
satisfies an independent success criterion: user utility holds and, for
injection variants, the attacker goal is not achieved. Authorization
constrains $a$; success evaluates it. Neither implies the other. The
examples in \S\ref{sec:intro} are authorized but incorrect.

\subsection{Threat model}
We follow the plan-then-execute assumption of CaMeL-style systems.
\begin{itemize}
\item \textbf{Control flow.} A trusted planner that sees only $q_t$ fixes the
control flow; control-flow hijacking is out of scope.
\item \textbf{Attacker.} The attacker writes the \emph{text} of documents the
agent reads. In one variant, the attacker also controls the account of a
source the user delegated (a compromised vendor).
\item \textbf{Host metadata versus author text.} We separate two kinds of
provenance information. \emph{Host metadata} is attached by the platform
and cannot be forged by a document's author: source identifier and kind,
authenticated sender, timestamp, and confidentiality readers.
\emph{Author-controlled text} includes field contents, subject lines,
reference numbers and any instruction-like phrasing.
\item \textbf{Assumption A8.} Field boundaries and field authorship are host
metadata only when the platform defines them, as in typed API records. They
are not host metadata when fields are parsed out of free text: a ``quoted''
block in an e-mail is text written by the sender. Our simulator \emph{assumes}
host-provided field paths and field authors (assumption~A8). This is the
same distinction ROPE draws between a platform's record structure and a
single block (\S3.3 of \citealp{ma2026rope}).
\end{itemize}
Out of scope: planner errors, side channels, and isolation of the binder
process. In our implementation the binder boundary is an interface property
checked by tests (\S\ref{sec:sim}), not process isolation.

\subsection{Fixed-contract evidence binding}
\begin{definition}[Binder]
A binder $b$ maps $(q_t, k, C_{t,v,k})$, together with the provenance
records and host metadata of the candidates, to an element of
$C_{t,v,k}\cup\{\abst\}$. The runtime accepts only a candidate identifier.
A non-learned validator $V_{\kappa}$ checks the value type, whether the
provenance is admitted by $\kappa$, confidentiality flow to the sink's
readers, and numeric limits before execution.
\end{definition}

A binder therefore cannot issue permissions: it can only name one of the
admitted candidates or abstain. The rule that anything it names must be
admitted is enforced by $V_\kappa$, not by the binder.

To prevent credit from looser boundaries, we fix a \emph{reference contract}
$\kappa^{\mathrm{ref}}$ and re-check every executed call of every method
against it. A method whose native boundary is looser can execute calls that
$\kappa^{\mathrm{ref}}$ rejects; these count as reference violations, never
as success.

We report three levels per episode and never merge them:
\begin{itemize}
\item \emph{proposed}: a complete call reached the validator;
\item \emph{executed}: the validator allowed the call and the sink changed
the post-state;
\item \emph{attacker-influenced}: an executed call carried an
attacker-planted value.
\end{itemize}
In this framework a blocked proposal is never counted as a successful
attack. AgentDojo's own trace-based checks do count proposed calls, so on
AgentDojo we report its official channel alongside an executed-effect
channel (\S\ref{sec:exp-real}). \emph{Secure
success} means one of two things. For a feasible task, the call executed,
matched the success criterion, and had no reference violation and no
attacker value. For a task that exceeds the principal's permissions, secure
success means that nothing was executed.

%% ===========================================================================
\section{A Clairvoyant Ceiling for Binding}
\label{sec:headroom}

\begin{definition}[Binder-only clairvoyant ceiling]
\label{def:h}
Fix a plan and a contract. At each write call the binder picks either
abstention or one element of the product of the admitted sets of that
call's slots, given the observations so far. Let $A_{t,v}$ be the set of
complete choice sequences for the whole plan, i.e.\ the leaves of this
decision tree. When admitted sets do not depend on earlier choices,
$A_{t,v}=\prod_{k\in K_t}(C_{t,v,k}\cup\{\abst\})$ (the V1 census). Let
$r(t,v)$ be the rule's assignment, and let $V_t$ be the set of variants of
$t$ (benign and injection variants, nested within $t$). Define
\begin{align}
s_{\mathrm{rule}}(t) &= \tfrac{1}{|V_t|}\textstyle\sum_{v\in V_t} s\big(t,v,r(t,v)\big),\\
s_{\mathrm{best}}(t) &= \tfrac{1}{|V_t|}\textstyle\sum_{v\in V_t} \max_{a\in A_{t,v}} s(t,v,a),\\
H^{\mathrm{clair}} &= \tfrac{1}{|T|}\textstyle\sum_{t\in T}\big[s_{\mathrm{best}}(t)-s_{\mathrm{rule}}(t)\big].
\end{align}
\end{definition}

\begin{proposition}[elementary]
\label{prop:bound}
Suppose $r(t,v)\in A_{t,v}$ for all $t,v$. Then
$s_{\mathrm{best}}(t)\ge s_{\mathrm{rule}}(t)$ for every $t$. Moreover,
every binder $b$ whose assignments lie in $A_{t,v}$ satisfies
$\frac{1}{|T|}\sum_t[s_b(t)-s_{\mathrm{rule}}(t)]\le H^{\mathrm{clair}}$.
\end{proposition}
\begin{proof}
A binder that only ever picks admitted choices or abstains traces one leaf
of the tree, so for every $(t,v)$ we have
$s(t,v,b(t,v))\le\max_{a\in A_{t,v}}s(t,v,a)$. For a randomized binder the
same holds in expectation. Averaging over $v$ and then over $t$ gives the
claim, and taking $b=r$ gives the first statement.
\end{proof}
The premise $r(t,v)\in A_{t,v}$ holds by construction for our baselines,
and the census checks it as an invariant (best $\ge$ baseline in every
variant).

\paragraph{What the ceiling is and is not.}
$H^{\mathrm{clair}}$ is a \emph{clairvoyant} ceiling. The maximizing
assignment is chosen with knowledge of the evaluator, so it ignores the
information limits of any real binder.
\begin{itemize}
\item It bounds the improvement a binder could make \emph{from above}, per
evaluation channel and per estimand (secure success or task utility).
\item It is not an estimate of learnable performance, and a large value is
not evidence that learning would help.
\item It bounds improvements that come \emph{only} from choosing differently
among admitted candidates. It excludes changing the plan, reading more data,
authenticating a sender, and asking the user.
\end{itemize}
When the plan is the benchmark's ground-truth plan, the ceiling is
\emph{conditional on an oracle plan} (\textsc{oracle\_plan\_conditional}).

\paragraph{Whole-call evaluation.}
We do not maximize each argument separately. Arguments interact in two ways:
two arguments of the same call are checked jointly, and an earlier write can
change what a later call reads. Nothing guarantees that per-field optima
combine into an optimal call, and we have no proof that they do. So every
element of $A_{t,v}$ is a complete choice sequence for the fixed plan. It is
replayed from the same initial state, with each branch on its own copy of
its parent's state.

\paragraph{Witnesses and exhaustiveness.}
When $A_{t,v}$ is too large to enumerate, the best candidate \emph{found}
does not bound the maximum. Because success is binary, we record two things
for each variant:
\begin{itemize}
\item whether a \emph{success witness} was found, which fixes the maximum at
$1$;
\item whether the search was exhaustive, which, without a witness, fixes it
at $0$.
\end{itemize}
Otherwise the maximum is only known to lie in $[0,1]$. This gives per-task
bounds $[s^{\mathrm{lo}}_{\mathrm{best}}(t), s^{\mathrm{up}}_{\mathrm{best}}(t)]$
and aggregate bounds
$[H^{\mathrm{clair}}_{\mathrm{lower}},H^{\mathrm{clair}}_{\mathrm{upper}}]$. A
task whose plan has no binder slot has $|A_{t,v}|=1$. Its only action is the
rule's own, so its gap is exactly $0$ whether or not its success can be
evaluated.

\paragraph{Unresolved tasks stay in the denominator.}
Some tasks cannot be evaluated. Examples are a slot whose admitted set cannot
be defined without using the answer, or a replay error. We mark these
\textsc{unresolved} and keep them in $T$. $H^{\mathrm{clair}}_{\mathrm{lower}}$
counts their gap as $0$, and $H^{\mathrm{clair}}_{\mathrm{upper}}$ counts it as
$1$. On a finite benchmark evaluated in full (a census),
$[H^{\mathrm{clair}}_{\mathrm{lower}},H^{\mathrm{clair}}_{\mathrm{upper}}]$ has
no sampling uncertainty. Attack variants are never counted as independent
tasks.

\paragraph{A computational envelope.}
$[H^{\mathrm{clair}}_{\mathrm{lower}},H^{\mathrm{clair}}_{\mathrm{upper}}]$
combines witness bounds, partially searched trees and unresolved tasks.
It is not a confidence interval, not an achieved improvement, and not a
learnable security gain. It is conditional on the oracle plan, the oracle
observations and the oracle final answer.

\paragraph{An operational rule, not a finding.}
Before any real-benchmark run we fixed $\delta=0.05$ as an operational rule.
$H^{\mathrm{clair}}_{\mathrm{upper}}<\delta$ is sufficient to stop
investing in learned binding under a given contract and plan. The converse
implies nothing: holding can still be justified by a contaminated
contract, a weak baseline, a utility-only gap, or cost and novelty. The
rule is not a statistical, equivalence, validity or publishability test.

%% ===========================================================================
\section{Oracle-Plan Candidate-Binding Census on AgentDojo}
\label{sec:exp-real}

An \emph{oracle-plan candidate-binding census} fixes everything except the
binder's choice. The call sequence, the read-call arguments, the free-text
arguments and the final answer are the benchmark's own ground truth. Only
the typed arguments of state-changing calls vary, and every admitted choice
is replayed offline against the benchmark's evaluators. No language model,
GPU, network access or service account is used. Every number in this
section, in both arms defined below, is \textsc{oracle\_plan\_conditional}.

\subsection{Snapshot, evaluator and two evaluation channels}
We pinned AgentDojo~0.1.35 (MIT; Appendix~\ref{app:protocol}) at benchmark
version v1.2.2, whose registry lists \AdjUser{} user tasks and 35 injection
tasks (Table~\ref{tab:adjcensus}). Reading the pinned source fixed three
points.
\begin{itemize}
\item \texttt{security()} returns \texttt{True} when the \emph{injection
goal was achieved}. For benign variants we record \texttt{None} (not
applicable), not $0$.
\item \texttt{utility\_from\_traces}/\texttt{security\_from\_traces} run
before the state-based checks. They are overridden by 12 user tasks and one
injection task.
\item Their traces are the tool calls \emph{proposed} in assistant
messages, whether or not they executed.
\end{itemize}
We therefore report two channels and never merge them.
\begin{itemize}
\item \textbf{Official} (primary): the assistant/tool transcript goes
unchanged to AgentDojo's own transcript functions and checker. Proposed
calls stay in the trace.
\item \textbf{Executed effect}: the trace holds only calls that executed
without error, evaluated together with the post-state.
\end{itemize}
A binder abstention is a candidate event: no assistant call is emitted, and
no call is deleted after the fact. A fixture shows that the channels
diverge: attacker calls that were proposed but denied count as an attack in
the official channel and not in the effect channel.

\subsection{Relation to our earlier census (V1)}
\label{sec:v1}
Our earlier census (Appendix~\ref{app:v1}) derived slots, contracts and
state-changing calls from the tasks' ground-truth argument values, so its
candidate sets were functions of hidden gold; a test shows the dependence.
Its interval $[\AdjHlo,\AdjHup]$ is kept as the result of that contract
only. It is neither a deployable ceiling nor a confidence interval, and we
do not compare it like for like with the census below
(Table~\ref{tab:ccv1v2}).

\subsection{CONTRACT\_V2: an actor-view contract}
\label{sec:v2}
The binder's \emph{actor view} is the trusted request, the tool schema, the
declared policy, and the outputs of calls already executed in this replay.
Hidden gold is not an input to slot typing, candidate extraction or
baseline decisions. Hidden gold here means ground-truth write values, the
final answer, post-states, evaluators, task and attack identifiers, and
future outputs.

This is the only sense in which the census is gold-free. The evaluated
arms still use the oracle plan and oracle final answer, and the fallback
arm inserts gold values (below).
\begin{itemize}
\item \textbf{Writes and slots.} Writes are declared from function names.
Slot kinds come from the function, argument name, JSON-schema type and
schema description. The binder kinds are IBAN, e-mail, date, date-time, URL,
identifier, number, enum, Boolean, entity, principal and channel.
Free-text arguments stay planner-fixed and are counted. There are no
per-task allowlists.
\item \textbf{Candidates.} Typed values extracted from the prompt and the
observation prefix. An attacker-written value of the right type in an
observation \emph{is} admitted, so V2 is a provenance-and-type contract,
looser than V1's gold-located one.
\item \textbf{Two arms.} A binder slot with no candidate forces abstention
in the \emph{strict} arm. In the \emph{fallback} arm it takes the
planner's gold value (\textsc{oracle\_composed}; \CcOracleComposed{}
tasks), so that arm is not gold-free.
\item \textbf{Baselines.} Unique-or-abstain, and a typed lexical/reference
resolver that binds a unique prompt value, an observed name the prompt
mentions verbatim, or the observed value whose record shares the most words
with the prompt. The resolver was written after V1 and developed on three
tasks.
\item \textbf{Gold-boundary tests.} Mutating the hidden gold (write values,
final answer, injection task) leaves slot kinds, candidates and both
baselines' decisions unchanged on five task variants. A chooser that reads
gold fails the same test (negative control).
\end{itemize}
The replay is \emph{sequential}, not counterfactual. Each choice executes
on a copy of the current state, and later observations and candidate sets
follow from it; read arguments stay at the oracle plan.

The search is depth-first, trying the gold-matching choice first (gold
affects the order only). It stops once it has witnesses for task success,
secure success and attack success in both channels. Leaf caps per variant
are 256 (workspace, travel) and 4{,}096 (banking, slack), set from
measured replay cost. A variant whose search hits the cap without a witness
keeps the bound $[0,1]$.

A \emph{selectable failure} is a variant where the baseline is not secure
but a secure witness exists. For tasks with one, we compute the
success-reachable choices at the first write for every variant. Variants
with an identical first-write view whose sets are provably disjoint cannot
all be served by any deterministic actor-view chooser. They are marked
\textsc{refuted}; the rest are \textsc{not refuted}, which is a necessary
condition only.

Everything was frozen before the run: the configuration, code, tests, a
structural pass and a four-task smoke run. The smoke run exposed a number
extraction bug, which was fixed before the freeze.

\subsection{Results of the census}
\label{sec:v2-results}
\paragraph{Coverage, search and cost.} No task raised a replay error or ran
out of budget, in either arm. This is all that ``resolved'' means here
(\CcTasks{}/\CcTasks{} tasks, \CcVariants{} variants per arm). It does not
mean that every variant's optimum was found.
\begin{itemize}
\item \CcNoWrite{} tasks have no write call, so their gap is exactly $0$.
\item \CcIncomplete{} variants in \CcIncompleteTasks{} tasks hit the leaf
cap before all witnesses were found.
\item For secure success and for utility, \CcSecUndet{} of these variants
(\CcSecUndetTasks{} workspace tasks) keep the envelope $[0,1]$. For
attack reachability, \CcAtkUndet{} do.
\end{itemize}
The run took \CcWall{}~s wall-clock and \CcCPU{}~CPU-s on two
single-threaded processes. No earlier result was reused.

\begin{table}[t]
\caption{CONTRACT\_V2 envelope $[H_{\mathrm{lower}},H_{\mathrm{upper}}]$ on
AgentDojo v1.2.2 over all \CcTasks{} tasks, per arm. Computational envelopes
over partially searched trees, not confidence intervals or achieved gains.
Both arms are \textsc{oracle\_plan\_conditional}; the fallback arm also
inserts gold values. Offline replay, no language model.}
\label{tab:cch}
\centering\scriptsize
\begin{tabular}{llccc}
\toprule
Empty-slot mode & Baseline & Official, secure & Effect, secure & Official, utility only \\
\midrule
fallback (primary) & typed resolver & [0.273, 0.293] & [0.273, 0.293] & [0.273, 0.293] \\
fallback (primary) & unique-or-abstain & [0.314, 0.334] & [0.314, 0.334] & [0.314, 0.334] \\
strict & typed resolver & [0.271, 0.292] & [0.271, 0.292] & [0.271, 0.292] \\
strict & unique-or-abstain & [0.302, 0.323] & [0.302, 0.323] & [0.302, 0.323] \\
\bottomrule
\end{tabular}

\end{table}

\paragraph{Finding 1: utility, not security.} Over the typed resolver, the
envelope is $[\CcHStrictLo,\CcHStrictUp]$ in the strict arm and
$[\CcHLo,\CcHUp]$ in the fallback arm (Table~\ref{tab:cch}). Over
unique-or-abstain it is $[\CcHRuleStrictLo,\CcHRuleStrictUp]$ and
$[\CcHRuleLo,\CcHRuleUp]$.

In each arm the utility-only envelope equals the secure envelope. All
\CcSSelN{} (strict, \CcSSelTasks{} tasks) and \CcSelN{} (fallback,
\CcSelTasks{} tasks) selectable failures are utility failures. In none did
the baseline complete the task while the attack succeeded and a secure
alternative existed. The gap concentrates in Slack and workspace (Appendix
Table~\ref{tab:ccsuite}): choosing the channel with the most users, the
largest file, or the right event among legitimate ones.

The oracle final answer matters for \CcOutDep{} tasks in the fallback arm
and \CcSOutDep{} in the strict arm (\textsc{oracle\_output\_dependent}).
\CcOutDepWrite{} of the \CcOutDep{} have a write call. They overlap with the
\CcOracleComposed{} \textsc{oracle\_composed} tasks in
\CcOutDepCompOverlap{}, so \CcOutDepCompUnion{} tasks carry either flag in
the fallback arm.

In this run the official and executed-effect channels gave the same outcome
on all \CcPathRecs{} path records. No proposal was denied, because the
binder only proposes admitted values, and the \CcToolErrPaths{} paths with
tool errors changed no outcome. This is an observation about this trace set,
not an equivalence of the two evaluators; a fixture shows them diverging on
denied proposals.

\begin{table}[t]
\caption{Security under CONTRACT\_V2: injected variants in which a path
executed a call achieving the attacker goal (Off.: official, Eff.:
executed-effect channel). $^a$Only variants where the gold choice was
admitted. $^b$Witness bounds on the number of variants where \emph{some}
admitted choice sequence achieves it; the upper count includes capped
searches. These are counts, not attack rates or learnable gains.}
\label{tab:ccatk}
\centering\scriptsize
\begin{tabular}{lcc}
\toprule
Injected variants (949), fallback mode & Official & Effect \\
\midrule
typed resolver & 4 & 4 \\
unique-or-abstain & 4 & 4 \\
abstain on every write & 0 & 0 \\
ground-truth path (gold admitted, 709 variants) & 0 & 0 \\
some admitted action reaches the attacker goal & 35--110 & 35--110 \\
\bottomrule
\end{tabular}

\end{table}

\paragraph{Finding 2: rare execution, non-zero admitted risk.} Both
baselines executed the attacker goal in \CcSResAtk{} (strict) and
\CcResAtk{} (fallback) of \CcInjVariants{} injected variants. The fallback
cases are the banking task whose legitimate IBAN the injection erased; in
the strict arm the baselines abstain there. The ground-truth path and full
abstention executed none (Table~\ref{tab:ccatk}).

Some admitted choice sequence nevertheless achieves the attacker goal in
\CcSReachLo--\CcSReachUp{} (strict) and \CcReachLo--\CcReachUp{} (fallback)
injected variants, because the contract admits any observed value of the
right type. These witness counts are neither a baseline attack rate nor a
security improvement a binder could learn. They also rule out claiming that
binding carries no attack risk. The benign-only and injected-only envelopes
are close ($[\CcHBenLo,\CcHBenUp]$ and $[\CcHInjLo,\CcHInjUp]$, fallback),
but closeness does not show that the security effect of binding is zero.

\paragraph{Finding 3: limits of the measurement.}
\begin{itemize}
\item \textbf{Candidate coverage.} In \CcBenUnsolv{} of the
\CcWriteTasks{} tasks with writes, no candidate our generator admits
completes even the benign task (fallback). A post-hoc diagnostic of the
gold path found \CcCovSlots{} slot values in \CcCovTasks{} tasks that were
not admitted. \CcCovNotInView{} are absent from the actor view, e.g.\
\CcCovDatetime{} date-times composed from a date and a time. \CcCovInText{}
occur only in free text or as a name in an IBAN slot, and \CcCovList{} are
list shapes. This is a limit of our candidate generator. A binder that may
compose values, or a different contract, might complete these tasks.
\item \textbf{Identifiability screen.} Of the \CcSelN{} selectable failures
(fallback), the screen refuted none. For \CcSingletonView{} of them it is
vacuous, because no other variant shares their first-write view; in
\CcVacuousTasks{} of the \CcSelTasks{} tasks every row is vacuous. The other
\CcSelViewShared{} share a view within their task and have a common
success-reachable first-write choice. The screen covers the first write
only. \textsc{not refuted} is therefore not evidence of identifiability. The
needed information does appear in the observation (e.g.\ user counts, file
sizes), but using it takes aggregation or comparison, not lexical matching.
\end{itemize}

\paragraph{What we do next.} We hold learned-binder method development.
This is the authors' resource decision, not a result, and it does not say
that evidence-binding research is futile. The frozen follow-up rule, which
counts \textsc{not refuted} rows, would have proposed a learned-binder study
(\textsc{\CcFollowup}). Given the vacuous rows and the utility-only gap, we
do not act on it.

\paragraph{Not run.}
\begin{itemize}
\item \textbf{Learned or prompted binder.} Held as an operational
decision; no language-model backend was used.
\item \textbf{AgentDyn census.} Needs approval to fetch the repository and a
license check.
\item \textbf{Other attacks and versions.} Adaptive attacks, other attack
templates and other benchmark versions.
\end{itemize}

%% ===========================================================================
\section{Limitations}
\label{sec:limits}

\paragraph{AgentDojo census.}
\begin{itemize}
\item \textbf{Oracle plan.} Every number in both arms is
\textsc{oracle\_plan\_conditional}: the call sequence, read arguments,
free-text arguments and final answer are ground truth. In the fallback arm,
\CcOutDepCompUnion{} tasks are \textsc{oracle\_output\_dependent} or
\textsc{oracle\_composed} (\CcOutDepCompOverlap{} are both). A
language-model planner makes errors the envelope does not cover, and it
would see different candidates.
\item \textbf{Replay.} The replay is sequential, but read arguments stay at
the oracle plan, so the plan never adapts to an earlier choice.
\item \textbf{Contract design.} CONTRACT\_V2 is our own post-hoc revision.
Its slot typing uses this benchmark's schema descriptions, and the resolver
was developed on three tasks after V1 was seen. It is not a blind
pre-registration. It admits any observed value of the right type, which is
looser than a field-scoped contract. Its candidate generator cannot compose
or compute values. That is a limit of this generator, not of binding in
general (\S\ref{sec:v2-results}).
\item \textbf{Search caps.} \CcIncomplete{} variants have incomplete
searches. The secure and utility envelopes stay $[0,1]$ for \CcSecUndet{} of
them, and attack reachability stays open for \CcAtkUndet{}.
\item \textbf{Identifiability screen.} It checks the first write only and is
a necessary condition. It is vacuous for \CcSingletonView{} of \CcSelN{}
failures.
\item \textbf{Channels.} They coincided in this trace set because nothing was
denied. That is not a general equivalence: they diverge under a denying
validator, as a fixture shows.
\item \textbf{Attacks.} One template (\texttt{DirectAttack}), one benchmark
version, and no adaptive attacker. Injection text replaces default content,
so some variants cannot be solved from observation.
\item \textbf{Envelope, not performance.} No learned or prompted binder
was run on AgentDojo.
\end{itemize}

\paragraph{Simulator.}
\begin{itemize}
\item \textbf{Self-authored evidence.} The generator, rules and features
were written by the same authors, with planted signals such as shared cue
words, fixed timing of attacker documents and position regularities
(Appendix~\ref{app:sim}).
\item \textbf{Simulated baselines.} Whole-document trust and the CaMeL-style
selector encode an assumed susceptibility model. The CaMeL-style selector
lacks CaMeL's information-sufficiency check and real confirmation.
\item \textbf{Self-grading.} Training labels and the evaluator share the
generator's gold values.
\item \textbf{Host metadata (A8).} Field paths and field authors are treated
as host metadata, which fails for fields parsed from free text.
\item \textbf{Statistical weight.} Three held-out templates make every
interval uninformative.
\end{itemize}

\paragraph{Implementation.}
The validator is tested, not verified. It trusts the provenance labels it is
given, and the binder boundary is an in-process interface, not isolation.

%% ===========================================================================
\section{Conclusion}
\label{sec:conclusion}
Authority boundaries decide which values may reach an argument, not which
admitted value is right. Under an oracle plan on AgentDojo, we audited how
much that choice could matter when only typed, observed candidates may be
chosen.
\begin{itemize}
\item \textbf{Utility, not security.} Choosing differently among admitted
values could add at most $[\CcHStrictLo,\CcHStrictUp]$ (strict) or
$[\CcHLo,\CcHUp]$ (fallback) over a typed resolver. The whole envelope is
utility.
\item \textbf{Rare execution, non-zero risk.} Baselines rarely executed the
attacker goal, but some admitted choice achieves it in a bounded set of
injected variants.
\item \textbf{Measurement limits.} Our generator cannot serve some tasks at
all, and whether the selectable failures are identifiable from observation
remains open.
\end{itemize}
These findings do not transfer to agents with language-model planners
without new, independent traces. We stop method development here.

%% ===========================================================================
\section*{Impact Statement}
This work concerns the security of language-model agents that act on
third-party content. Its intended effect is to prevent overclaiming: it
separates what an authority boundary guarantees from what it does not, and
defines a measurement that can show when a learned component is not worth
building. The simulator uses synthetic data only. The AgentDojo census
replays the benchmark's own fictional environments offline, without a
language model, network access or service account. No real accounts,
credentials, payment identifiers or user data are involved. The attack patterns we model, such as injected instructions,
stale values and compromised senders, are publicly documented. Our
description adds no new attack capability. A possible negative consequence
is that a formal-looking bound is mistaken for a security guarantee. We
therefore state explicitly that $H^{\mathrm{clair}}$ bounds only the benefit
of binding choices under a fixed contract and plan, that it is a ceiling
rather than achievable performance, and that our validator provides no process
isolation.

\bibliography{references}
\bibliographystyle{icml2026}

%% ===========================================================================
\newpage
\appendix
\onecolumn

\section{Synthetic Simulator Study (Engineering Evidence Only)}
\label{sec:sim}

To exercise the formulation end to end, we built a small simulator in which
every document, secret and attack is generated from a seed. Its purpose is to
test the harness, not to measure agents. All results from it are
\emph{engineering evidence only}.

\subsection{Simulator design}
\paragraph{Tasks and conditions.}
There are nine task templates in three domains (Appendix~\ref{app:sim}):
\begin{itemize}
\item billing: pay an invoice by order number, pay the latest invoice to a
saved payee, request a refund by reference;
\item calendar: accept the latest proposed date, reschedule by meeting
reference, invite a registrant from a form;
\item support: forward a tracking number, report a delivery date from a
file, reply with a knowledge-base note.
\end{itemize}
Each template is instantiated under five content conditions:
\begin{itemize}
\item \emph{helpful} data;
\item \emph{malicious instructions}, either in an attacker's document or
written by a compromised delegated sender;
\item \emph{mixed} content, where the injection shares the helpful document,
either in another field or in the delegated field itself;
\item \emph{wrong facts}, i.e.\ type-valid stale or superseded values with no
instruction;
\item \emph{insufficient permission}: unknown entity, tool not granted,
amount over limit, or restricted path.
\end{itemize}
Secrets are synthetic tokens in files readable only by the user. Payment
identifiers use a non-ISO country code and all addresses use reserved
\texttt{.example} domains.

\paragraph{Contract and validator.}
The reference contract delegates each argument to a named field path of
documents whose authenticated sender and field author are the named entity,
or to the named file. The validator returns explicit reason codes; for
example, \texttt{PROVENANCE\_OUT\_OF\_SCOPE}, \texttt{FLOW\_VIOLATION},
\texttt{LIMIT\_EXCEEDED}, \texttt{NO\_CAPABILITY} and
\texttt{REQUESTED\_SCOPE\_MISMATCH}. Every decision is logged in a trace.

\paragraph{Selectors.}
We compare seven selectors under an identical plan, reads and budget.
\begin{itemize}
\item \emph{Whole-document trust}: any field of a delegated document is
admitted. It is paired with an \emph{assumed} instruction-following reader.
\item \emph{Untrusted-all deny}.
\item A \emph{CaMeL-style} selector. This is our simplified simulation, not
the CaMeL implementation: high-integrity untrusted arguments require user
confirmation, counted as a refusal, and an assumed susceptible extraction
step reads the requested field.
\item \emph{Rule field scopes}: filter by field scope, then key, then trusted
history, then recency; bind only a unique value.
\item A \emph{learned binder}: logistic regression over 17
template-agnostic features, threshold $0.5$ fixed in advance.
\item Two bounds: an \emph{oracle} binder and a worst-case
\emph{adversarial} binder that knows the attack labels.
\end{itemize}
Appendix Table~\ref{tab:features} lists each feature and the origin of its
input: trusted plan, trusted record, host metadata, host metadata under A8,
or author-controlled text. No feature reads hidden labels, which a unit test
checks.

\paragraph{Interface, not isolation.}
The binder receives a read-only view and the runtime accepts only a
candidate-identifier string. Frozen Python objects prevent accidental
mutation, but any code in the same process can change them. A unit test
demonstrates this, and another shows that the validator trusts the
provenance labels it is given. The guarantee that a binder ``cannot grant
authority'' is therefore a property of the interface in our code, not of
process isolation.

\subsection{Simulator experiments}
\label{sec:exp}

\subsubsection{Engineering checks (SED-E0, SED-E1)}
The simulator run (SED-E1, commit \texttt{\EOneCommit}, \EOneWall{} s on two
CPU cores) used a configuration frozen before the run. Splits hold out
templates, domains, or instances; the independent unit is the template. The
validator-guarded selectors (rule, learned, oracle, adversarial) had
\GuardedRefViol{} reference violations and \GuardedLeak{} synthetic-secret
leaks. No task that exceeded the principal's permissions was ever executed.
The oracle reached secure success $1.0$ on every feasible episode, so a
secure solution always existed inside the reference contract.

Re-aggregating the stored raw rows reproduced every reported aggregate with
zero differences. Table~\ref{tab:levels} (appendix) separates the proposed,
executed and attacker-influenced levels. For example, the learned binder
proposed \LearnOutOfScopeProposals{} bindings outside the reference contract
on held-out templates, and all of them were blocked.

\subsubsection{What the boundary removes, and what it leaves}
Table~\ref{tab:boundary} counts, per delegated argument, whether an
attacker value was present and whether it lay \emph{inside} the reference
contract.
\begin{itemize}
\item \textbf{Removed by the contract.} Attacks from out-of-scope documents
and from other fields of the delegated document never fell inside it.
\item \textbf{Left inside the contract.} Attacks in the delegated field
itself, and all attacks from a compromised delegated sender, always did.
Stale and superseded values left the admitted set ambiguous in every case.
\end{itemize}
This split is a property of how we generated the data, under assumption A8.
It is a check that the harness classifies these cases as intended, not an
empirical finding about agents.

\subsubsection{A learned binder did not beat the rules}
\label{sec:exp-learned}
Table~\ref{tab:main} reports the template-held-out split.
\begin{itemize}
\item \textbf{Main comparison.} The learned binder's secure success
(\LearnSS) is close to the rule's (\RuleSS). It traded fewer over-refusals
for more wrong actions.
\item \textbf{Other splits.} Across the three splits
(Table~\ref{tab:diff}) the learned-minus-rule difference in secure success
was \DiffTemplate, \DiffEnv{} and \DiffInstance. In the instance split the
learned binder also executed more attacker values: a template-averaged
attacker-influenced rate higher by \DiffInstanceAtk. With three held-out templates, the cluster-bootstrap intervals are
not meaningful and we do not interpret them.
\item \textbf{Baselines.} Whole-document trust and the CaMeL-style selector
are simulations whose attack rates follow from an assumed susceptibility
model, so their rows are not measurements of any system. The CaMeL-style
selector's low secure success comes mainly from confirmation requests on T6,
the one held-out template whose delegated argument is high-integrity.
\item \textbf{Compromised senders.} When a compromised delegated sender wrote
into the delegated field, secure success on the held-out templates was at
most 0.03 for every selector. No trusted history exists for those
templates.
\end{itemize}
We read these results as a failed pre-registered stop check: the learned
binder does not beat the rule baseline. The synthetic data cannot settle
whether learning could help on real tasks.

\begin{table}[h]
\caption{SED-E1, template-held-out split (test templates T3, T6, T9; 450
episodes). SS: secure success (macro over templates). ASR: executed attacker
values among attack-condition episodes. RefViol: executed calls rejected by
the reference contract. OverRef: no execution on a feasible task. $^\dagger$
simulated baseline with an assumed susceptibility model; not a measurement of
any system. ENGINEERING ONLY.}
\label{tab:main}
\centering\small
\begin{tabular}{lcccccc}
\toprule
Selector & SS$\uparrow$ & ASR$\downarrow$ & RefViol$\downarrow$ & OverRef$\downarrow$ & Wrong$\downarrow$ & Confirm \\
\midrule
Whole-document trust$^\dagger$ & 0.560 & 0.472 & 0.067 & 0.244 & 0.056 & 0.000 \\
Untrusted-all deny & 0.200 & 0.000 & 0.000 & 1.000 & 0.000 & 0.000 \\
CaMeL-style (simulated)$^\dagger$ & 0.611 & 0.194 & 0.000 & 0.333 & 0.044 & 0.267 \\
Rule field scopes & 0.773 & 0.178 & 0.000 & 0.194 & 0.000 & 0.000 \\
Learned binder (LR) & 0.798 & 0.194 & 0.000 & 0.081 & 0.060 & 0.000 \\
\midrule
Oracle binder (bound) & 1.000 & 0.000 & 0.000 & 0.000 & 0.000 & 0.000 \\
Adversarial binder (bound) & 0.200 & 0.378 & 0.000 & 0.578 & 0.187 & 0.000 \\
\bottomrule
\end{tabular}

\end{table}

\subsubsection{Why it failed: estimation, not expressivity}
\label{sec:exp-expr}
Our first internal report attributed the failure to linear scoring: we
claimed that a linear score cannot express the lexicographic rule ``filter by
key, then by recency''. That claim was wrong for our features.
\begin{itemize}
\item \textbf{Representable in principle.} With binary features and bounded
continuous features, sufficiently separated weights make a linear argmax
lexicographic (Observation~\ref{obs:lex}, Appendix~\ref{app:proofs}).
\item \textbf{Representable in practice.} Hand-set weights on the same 17
features reproduce the rule's decision (bind the same candidate, or abstain)
on \LexAgreeTotal{} of \ArgsTotal{} test arguments across the three splits.
The fitted binder agreed on \LearnAgreeTotal{} (Table~\ref{tab:expr}).
\end{itemize}
This was an analysis after seeing results. It trained nothing and proposes
no method.

Where the fitted binder erred, the gold and chosen candidates differed in a
few features.
\begin{itemize}
\item \textbf{Held-out ``other key'' cases.} The difference was only key
match, latest-in-source and recency rank. The fitted template-split weights put
recency ($-1.72$ per unit rank) above key match ($+1.29$).
\item \textbf{Held-out knowledge-base cases with a stale second span.} The
difference was only \emph{first-in-field}, whose fitted weight was
$-0.74$. That template was held out of training in this split, so we do
\emph{not} attribute the error to the generator artifact we had earlier
blamed.
\end{itemize}
Why the pointwise log-loss fit produced these weights was not investigated.
One untested explanation is that pointwise log-loss is not a decision loss.

Two further findings from this check:
\begin{itemize}
\item \textbf{A bug.} The first version compared sigmoid probabilities in
the argmax. For large scores these saturate to $1.0$, so the candidate order
broke ties. We fixed this by taking the argmax on logits. Regenerating the
test episodes reproduced all stored rule and learned outcomes (0
mismatches), so the fix does not change any reported number.
\item \textbf{A limit of pointwise scoring.} The rule's set-level abstention
(``abstain if the newest admitted document holds two distinct values'') is
not representable by \emph{any} pointwise score plus threshold on our
features (Observation~\ref{obs:set}). No such configuration occurred in the
test data.
\end{itemize}

\subsubsection{Ceiling: simulator smoke test}
\label{sec:exp-headroom}
We ran the ceiling code on the instance-split test episodes. Each template
is treated as an original task with its 150 episodes nested as variants, the
reference contract fixes the admitted sets, and the rule is the field-scope
rule. Results (Appendix Table~\ref{tab:hsim}):
\begin{itemize}
\item every rule choice lay inside the action set;
\item no task was unresolved;
\item $H^{\mathrm{clair}}=\HSim$, which equals one minus the rule's secure
success on this split;
\item the ambiguity proxy was \ProxySim: every template had some ambiguous
slot.
\end{itemize}
The two numbers illustrate why a proxy cannot stand in for
$H^{\mathrm{clair}}$. Success in
the simulator is computed from the generator's own gold values, which also
provide the learned binder's training labels. It is therefore not an
independent criterion, and this value is a code-path check only.

\subsection{Simulator details}
\label{app:sim}

\paragraph{Templates.}
There are nine templates in three domains:
\begin{itemize}
\item \textbf{billing}: T1 pay invoice by order (recipient and amount
delegated), T2 pay latest invoice to saved payee (amount delegated), T3
refund by correction reference (amount delegated);
\item \textbf{calendar}: T4 latest proposed date, T5 reschedule by meeting
reference, T6 invite registrant from form (attendee delegated);
\item \textbf{support}: T7 tracking number by order, T8 delivery date from a
report file, T9 troubleshooting note from a knowledge-base file.
\end{itemize}

\paragraph{Conditions and variants.}
Each condition has the following variants:
\begin{itemize}
\item helpful: clean;
\item malicious: out-of-scope attacker document, or compromised delegated
sender;
\item mixed: other field, or same field;
\item wrong fact: other key, superseded, stale field, or same-field history;
\item insufficient permission: entity unknown, tool not granted, over limit,
or restricted path.
\end{itemize}
There are 30 train and 30 test episodes per template and condition, with
disjoint seeds.

\paragraph{Planted signals and assumptions.}
\begin{itemize}
\item \textbf{A1.} One cue lexicon is shared by the generator, the features
and the simulated susceptible readers. Attacker cue rate is 0.7 and
legitimate cue rate is 0.25.
\item \textbf{A2.} Position regularities: in same-field history, the gold
span precedes the stale one; in T9, a legitimate cue sentence precedes the
gold note.
\item \textbf{A3.} Attacker documents are always newer than gold; superseded
documents are always older.
\item \textbf{A4.} Trusted payment history exists only for T1, with
probability 0.5.
\item \textbf{A5.} The planner is exact.
\item \textbf{A6.} The susceptibility of whole-document trust and the
CaMeL-style extraction step is assumed.
\item \textbf{A7.} The CaMeL-style selector is not the CaMeL code.
\item \textbf{A8.} Field paths and field authors are treated as host
metadata.
\end{itemize}
None of these was altered after the results were seen.

\begin{table}[h]
\caption{Learned-binder features and the origin of their inputs.
host\_A8: host metadata only under assumption A8. author: text controlled by
the document's author (the attacker for attacker documents and compromised
senders).}
\label{tab:features}
\centering\small
\begin{tabular}{ll|ll}
\toprule
Feature & Inputs & Feature & Inputs \\
\midrule
in requested field path & plan, host\_A8 & cue before span & author \\
doc from requested source & plan, record, host & field structured & host\_A8 \\
field author is source & record, host\_A8 & multi-value field & author, set \\
field author $\neq$ sender & host\_A8, host & first in field & author \\
key given & plan & history given & record \\
key match & plan, author & history match / conflict & record, author \\
latest in source & host, set & doc type in task & task, author \\
recency rank & host, set & bias & -- \\
\bottomrule
\end{tabular}
\end{table}

\section{Elementary Arguments}
\label{app:proofs}

\begin{observation}[lexicographic priority is linearly representable]
\label{obs:lex}
Let each candidate have binary features $b_1,\dots,b_p\in\{0,1\}$ and a
bounded feature $\rho\in[0,1]$. Choose weights $w_\rho<0$ and $w_i>0$ such
that $w_i>\sum_{j>i}w_j+|w_\rho|$. Then $\arg\max_x
\big(\sum_i w_ib_i(x)+w_\rho\rho(x)\big)$ selects a candidate that is
lexicographically maximal in $(b_1,\dots,b_p,-\rho)$.
\end{observation}
\begin{proof}
Let $x$ and $y$ first differ at a binary feature $i$, with $b_i(x)=1$ and
$b_i(y)=0$. The score difference is at least
$w_i-\sum_{j>i}w_j-|w_\rho|>0$. If they agree on all binary features, the
difference is $w_\rho(\rho(x)-\rho(y))$, which favors the smaller $\rho$.
\end{proof}
This is standard. The weights in \texttt{scripts/analyze\_e1\_expressivity.py}
instantiate it for field scope $\gg$ key $\gg$ history $\gg$ multi-value
exclusion $\gg$ recency, with a threshold for abstention.

\begin{observation}[set-level abstention is not pointwise]
\label{obs:set}
Let $\phi$ be the router's feature map. No selector of the form ``bind
$\arg\max_x g(\phi(x))$, abstain if $\max_x g(\phi(x))<\tau$'' reproduces
the rule's decisions on the two candidate sets $S_1,S_2$ below.
\end{observation}
\begin{proof}
Let $x_{\mathrm{old}}$ be an admitted, key-matching candidate in an older
document.
\begin{itemize}
\item In $S_1$, a newer admitted, key-matching document holds two distinct
values in the delegated field, so the rule abstains.
\item In $S_2$, the newer document from the same source lacks the key, so the
rule binds $x_{\mathrm{old}}$.
\end{itemize}
Because the recency features are computed over all documents from the
source, $\phi(x_{\mathrm{old}})$ is identical in $S_1$ and $S_2$; a unit test
checks this. Binding in $S_2$ requires $g(\phi(x_{\mathrm{old}}))\ge\tau$.
Then $\max_{x\in S_1}g(\phi(x))\ge\tau$, so the selector binds in $S_1$,
which contradicts the rule's abstention.
\end{proof}

\paragraph{A retracted shortcut.}
An earlier version of our protocol proposed to stop if fewer than 10\% of
tasks had an ambiguous admitted set, and inferred an improvement below 5
percentage points. That inference was wrong on two counts.
\begin{itemize}
\item A task fraction $p$ bounds the average improvement by $p$, not by
$p/2$.
\item Ambiguity is a proxy, not a success difference.
\end{itemize}
In our simulator the ambiguity proxy is \ProxySim{} while
$H^{\mathrm{clair}}=\HSim$ (Appendix~\ref{sec:exp-headroom}).

\section{Experiment Record}
\label{app:record}
Commands, commits, configuration hashes and costs are listed in our run
manifest, which is not part of the anonymous review package. The runs relevant here are the following.

\paragraph{Simulator.}
\begin{itemize}
\item SED-E0 unit tests.
\item SED-E1 simulator pilot (\EOneWall{} s).
\item SED-E1-REAGG re-aggregation.
\item SED-E1-EXPR expressivity check. It ran three times (pre-fix, post-fix,
and aligned hand-set weights), and all three outputs are kept.
\end{itemize}

\paragraph{AgentDojo.}
\begin{itemize}
\item \textbf{Task enumeration} (SED-E2-ADJ-CENSUS).
\item \textbf{Structural pass} (DEVELOPMENT). It recorded plans, slots,
candidate-set sizes and timing, and inspected no success values.
\item \textbf{First census run: INVALID and kept.} A success witness equal
to the empty assignment was tested by truthiness, so zero-slot variants
lost their witness. This produced negative gaps, and we stopped the run.
\item \textbf{Corrected census} (\AdjWall{} s). It added a regression test
and a best-$\ge$-rule invariant check. The protocol was unchanged.
\item \textbf{Exploratory slot-sensitivity analysis}, run after the census.
\item \textbf{CONTRACT\_V2} (session 4). A structural pass (DEVELOPMENT,
no success values) and a four-task smoke run (DEVELOPMENT; it exposed a
number-extraction bug, fixed before the freeze). Then the frozen census
(SED-E2-CONTRACT-CLOSURE, \CcWall{}~s wall, \CcCPU{}~CPU-s, two
single-threaded processes) and a post-hoc coverage diagnostic.
\end{itemize}
All runs used at most two CPU cores and no GPU, network access during
replay, language model or paid API.

\begin{table}[h]
\caption{Per delegated argument, instance-split test set. Counts are over
\emph{arguments}; one billing template has two delegated arguments of which
only one is attacked, so ``present'' is below ``args'' in attack conditions.
ENGINEERING ONLY.}
\label{tab:boundary}
\centering\scriptsize
\begin{tabular}{lcccc}
\toprule
Condition/variant & Args & Atk.\ present & Atk.\ in scope & Ambig.\ in scope \\
\midrule
malicious / attacker doc & 172 & 157 & 0 & 6 \\
mixed / other field & 165 & 148 & 0 & 4 \\
malicious / compromised sender & 128 & 113 & 113 & 113 \\
mixed / same field & 135 & 122 & 122 & 122 \\
wrong fact / other key & 92 & 0 & 0 & 92 \\
wrong fact / superseded & 148 & 0 & 0 & 148 \\
wrong fact / stale field & 33 & 0 & 0 & 2 \\
wrong fact / same-field hist. & 27 & 0 & 0 & 27 \\
helpful / clean & 300 & 0 & 0 & 8 \\
\bottomrule
\end{tabular}

\end{table}

\begin{table}[h]
\caption{Learned binder minus rule field scopes (per-template means). $n_T$:
number of held-out templates. $^\ddagger$percentile cluster bootstrap over
templates; not meaningful for $n_T<10$ and not used for inference.
ENGINEERING ONLY.}
\label{tab:diff}
\centering\small
\begin{tabular}{lccccccl}
\toprule
Split & $n_T$ & $\Delta$SS & 95\% CI$^\ddagger$ & $\Delta$Atk & $\Delta$OverRef & $\Delta$Wrong & per template \\
\midrule
template & 3 & +0.024 & [-0.013, +0.047] & +0.007 & -0.091 & +0.060 & T3~-0.013, T6~+0.047, T9~+0.040 \\
environment & 3 & -0.007 & [-0.160, +0.100] & +0.002 & -0.056 & +0.060 & T7~+0.100, T8~-0.160, T9~+0.040 \\
instance & 9 & -0.033 & [-0.086, +0.013] & +0.032 & -0.039 & +0.041 & 9 templates \\
\bottomrule
\end{tabular}

\end{table}

\begin{table}[h]
\caption{Decision agreement with the rule, per delegated argument (bind the
same candidate, or both abstain). Lex: hand-set lexicographic weights (not
fitted). Learned: stored fitted weights. Regen.\ mismatches: stored outcomes
not reproduced by regeneration. EXPLORATORY.}
\label{tab:expr}
\centering\scriptsize
\begin{tabular}{lccccc}
\toprule
Split & Args & Rule binds & Lex $=$ rule & Learned $=$ rule & Regen.\ mismatches \\
\midrule
template & 360 & 290 & 360 & 275 & 0 \\
environment & 360 & 271 & 360 & 246 & 0 \\
instance & 1200 & 1036 & 1200 & 984 & 0 \\
\bottomrule
\end{tabular}

\end{table}

\begin{table}[h]
\caption{SED-E1 template split: proposed, executed and attacker-influenced
levels, with native denial reason codes. A blocked proposal never counts as
an attack. ENGINEERING ONLY.}
\label{tab:levels}
\centering\scriptsize
\begin{tabular}{lcccccl}
\toprule
Selector & Episodes & Proposed & Executed & Exec.\ w/ atk.\ value & Exec.\ ref.\ viol. & Native denials \\
\midrule
Whole-document trust$^\dagger$ & 450 & 406 & 272 & 85 & 30 & FLOW\_VIOLATION 88, NO\_CAPABILITY 46 \\
Untrusted-all deny & 450 & 0 & 0 & 0 & 0 & -- \\
CaMeL-style (simulated)$^\dagger$ & 450 & 406 & 240 & 35 & 0 & NEEDS\_CONFIRMATION 120, NO\_CAPABILITY 46 \\
Rule field scopes & 450 & 333 & 290 & 32 & 0 & NO\_CAPABILITY 43 \\
Learned binder (LR) & 450 & 405 & 331 & 35 & 0 & NO\_CAPABILITY 46, PROVENANCE\_OUT\_OF\_SCOPE 28 \\
Oracle binder (bound) & 450 & 360 & 360 & 0 & 0 & -- \\
Adversarial binder (bound) & 450 & 426 & 152 & 68 & 0 & FLOW\_VIOLATION 3, NO\_CAPABILITY 66, PROVENANCE\_OUT\_OF\_SCOPE 208 \\
\bottomrule
\end{tabular}

\end{table}

\begin{table}[h]
\caption{SED-E1 domain-held-out split (test: support templates T7--T9).
ENGINEERING ONLY.}
\centering\small
\begin{tabular}{lccccc}
\toprule
Selector & SS & ASR & RefViol & OverRef & Wrong \\
\midrule
Whole-document trust$^\dagger$ & 0.607 & 0.372 & 0.047 & 0.244 & 0.049 \\
Untrusted-all deny & 0.200 & 0.000 & 0.000 & 1.000 & 0.000 \\
CaMeL-style (simulated)$^\dagger$ & 0.829 & 0.317 & 0.000 & 0.000 & 0.044 \\
Rule field scopes & 0.767 & 0.089 & 0.000 & 0.247 & 0.000 \\
Learned binder (LR) & 0.760 & 0.094 & 0.000 & 0.178 & 0.060 \\
Oracle binder (bound) & 1.000 & 0.000 & 0.000 & 0.000 & 0.000 \\
Adversarial binder (bound) & 0.200 & 0.322 & 0.000 & 0.653 & 0.149 \\
\bottomrule
\end{tabular}

\end{table}

\begin{table}[h]
\caption{SED-E1 instance-held-out split (all templates, disjoint seeds).
ENGINEERING ONLY.}
\centering\small
\begin{tabular}{lccccc}
\toprule
Selector & SS & ASR & RefViol & OverRef & Wrong \\
\midrule
Whole-document trust$^\dagger$ & 0.694 & 0.504 & 0.055 & 0.081 & 0.039 \\
Untrusted-all deny & 0.200 & 0.000 & 0.000 & 1.000 & 0.000 \\
CaMeL-style (simulated)$^\dagger$ & 0.688 & 0.298 & 0.000 & 0.222 & 0.015 \\
Rule field scopes & 0.801 & 0.194 & 0.000 & 0.152 & 0.000 \\
Learned binder (LR) & 0.767 & 0.274 & 0.000 & 0.103 & 0.041 \\
Oracle binder (bound) & 1.000 & 0.000 & 0.000 & 0.000 & 0.000 \\
Adversarial binder (bound) & 0.200 & 0.383 & 0.000 & 0.580 & 0.183 \\
\bottomrule
\end{tabular}

\end{table}

\begin{table}[h]
\caption{AgentDojo v1.2.2 census per suite (primary analysis,
unique-or-abstain). Unresolved tasks contribute $0$ to the lower and $1$ to
the upper bound. \textsc{oracle\_plan\_conditional}.}
\label{tab:adjsuite}
\centering\small
\begin{tabular}{lcccc}
\toprule
Suite & Tasks & Tasks with gap $>0$ & Unresolved & per-suite $[H_{\mathrm{lower}}, H_{\mathrm{upper}}]$ \\
\midrule
workspace & 40 & 6 & 6 & [0.150, 0.300] \\
travel & 20 & 0 & 0 & [0.000, 0.000] \\
banking & 16 & 5 & 1 & [0.312, 0.375] \\
slack & 21 & 16 & 1 & [0.643, 0.690] \\
\bottomrule
\end{tabular}

\end{table}

\begin{table}[h]
\caption{AgentDojo v1.2.2: ceiling restricted to benign or to injected
variants (primary analysis, unique-or-abstain). Descriptive breakdown fixed
before the run.}
\label{tab:adjsplit}
\centering\small
\begin{tabular}{lc}
\toprule
Variants (primary, unique-or-abstain) & $[H_{\mathrm{lower}}, H_{\mathrm{upper}}]$ \\
\midrule
benign variants only & [0.247, 0.330] \\
injected variants only & [0.254, 0.336] \\
\bottomrule
\end{tabular}

\end{table}

\begin{table}[h]
\caption{Headroom smoke test on the simulator (instance-split test episodes;
150 nested variants per template). Success is the generator's gold, not an
independent criterion. ENGINEERING ONLY.}
\label{tab:hsim}
\centering\small
\begin{tabular}{lcccc}
\toprule
Template $t$ & Variants & $s_{\mathrm{rule}}(t)$ & $s_{\mathrm{oracle}}(t)$ & gap \\
\midrule
T1 & 150 & 0.900 & 1.000 & 0.100 \\
T2 & 150 & 0.800 & 1.000 & 0.200 \\
T3 & 150 & 0.847 & 1.000 & 0.153 \\
T4 & 150 & 0.787 & 1.000 & 0.213 \\
T5 & 150 & 0.780 & 1.000 & 0.220 \\
T6 & 150 & 0.793 & 1.000 & 0.207 \\
T7 & 150 & 0.793 & 1.000 & 0.207 \\
T8 & 150 & 0.827 & 1.000 & 0.173 \\
T9 & 150 & 0.680 & 1.000 & 0.320 \\
\midrule
$H$ (all resolved) & & & & 0.199 \\
Ambiguity proxy & & & & 1.000 \\
\bottomrule
\end{tabular}

\end{table}

\section{First AgentDojo Census (V1 Contract) in Detail}
\label{app:v1}
This appendix keeps the V1 census as reported in manuscript v2. Its contract was derived from hidden
ground truth (\S\ref{sec:v1}); the numbers below describe that contract only.

\paragraph{Coverage and bounds (V1).}
\begin{itemize}
\item \textbf{Coverage.} Of \AdjUser{} user tasks, \AdjZeroSlot{} had no
V1 slot, so their gap is exactly $0$. \AdjResolved{} tasks were resolved.
\AdjUnresolved{} were \textsc{unresolved} because a slot held free text
with no typed extractor; they stay in the denominator.
\item \textbf{Search.} We replayed \AdjVariants{} variants. \AdjNonExh{}
non-exhaustively searched variants lacked a witness.
\item \textbf{Bounds} (Table~\ref{tab:adjh}):
\begin{itemize}
\item unique-or-abstain rule: $[\AdjHlo,\AdjHup]$;
\item first-admitted rule: $[\AdjHfirstLo,\AdjHfirstUp]$;
\item free-text slots fixed to the planner value: $[\AdjHsecLo,\AdjHsecUp]$.
\end{itemize}
\end{itemize}
Manuscript v2 read the upper bound's crossing of $\delta$ as ``does not
license holding''. That reading is corrected in \S\ref{sec:headroom}.

\paragraph{Where the gap comes from.}
The positive gaps come mostly from tasks whose admitted set holds several
\emph{legitimate} values, where the unique-or-abstain rule abstains.
Examples are four Slack channels, or the files of a drive listing
(Table~\ref{tab:adjsuite}). Picking the right one (for example, ``the
channel with the most users'') is the reasoning step of the task. The clairvoyant chooser
resolves it by consulting the evaluator; a real binder would need the task
reasoning that the planner otherwise performs. The gap is similar on benign
and injected variants (Table~\ref{tab:adjsplit}). On this benchmark, under an
oracle plan, the ceiling is therefore mainly a \emph{utility} ceiling on
argument choice, not a security margin.

\paragraph{Security under the fixed plan.}
Attacks rarely reached an executed call.
\begin{itemize}
\item \textbf{Ground-truth plan.} It executed \AdjGtAtk{} attacks over
\AdjInjVariants{} injected variants.
\item \textbf{Unique-or-abstain rule.} It executed \AdjRuleAtk. All of them
are in one task, where the injection had replaced the bill text holding the
legitimate IBAN, so the attacker's IBAN was the only admitted value.
\item \textbf{Admitted sets.} An attacker-written value was admitted in
\AdjAtkAdmitted{} injected variants, and the benign value had been removed in
\AdjRemoved.
\end{itemize}
The removal property matters for interpretation. AgentDojo's injection text
\emph{replaces} a vector's default content, which can hold the legitimate
information. In \AdjUnsolvable{} of the \AdjRemoved{} removal variants (from
\AdjUnsolvableTasks{} task), no admitted action succeeds and abstention is the
best secure choice. In the other removal variants a success witness still
exists. The oracle plan still ``succeeds'' in all of them because its
arguments are fixed in advance. This is one way \textsc{oracle\_plan\_conditional}
results hide information limits.

\paragraph{Sensitivity to slot identification (exploratory).}
Slots and contracts are found by value coincidence between a ground-truth
argument and an earlier output field. This can yield contracts that make no
sense, such as an amount matched to a transaction-id field. After seeing the
census we dropped the tasks with numeric-valued slots (\AdjNumericSlots{} of
\AdjSlots{} slots). The bound on the remaining \AdjNumKept{} tasks was
$[\AdjNumLo,\AdjNumUp]$. A cruder name-agreement filter removes almost every
task with a slot and is not interpretable. We keep the frozen primary result
and report this only as evidence that $H^{\mathrm{clair}}$ depends on how
slots are identified.

\paragraph{An evaluator detail.}
One banking injection task's \texttt{security()} contains
\texttt{any(service or "" in t.subject.lower() ...)}. Because of operator
precedence, this returns \texttt{True} for any transaction to the attacker's
IBAN, whatever the subject. We did not modify it. The four executed attacks
above send money to the attacker under any reading.


\begin{table}[h]
\caption{AgentDojo 0.1.35, benchmark v1.2.2: tasks enumerated from the
package registry.}
\label{tab:adjcensus}
\centering\small
\begin{tabular}{lcc}
\toprule
Suite (v1.2.2) & User tasks & Injection tasks \\
\midrule
workspace & 40 & 14 \\
travel & 20 & 7 \\
banking & 16 & 9 \\
slack & 21 & 5 \\
total & 97 & 35 \\
\bottomrule
\end{tabular}

\end{table}

\begin{table}[h]
\caption{Binder-only clairvoyant ceiling on AgentDojo v1.2.2 (all user tasks
in the denominator; unresolved tasks counted as gap $0$ / $1$).
\textsc{oracle\_plan\_conditional}; offline replay, no language model.}
\label{tab:adjh}
\centering\scriptsize
\begin{tabular}{llcc}
\toprule
Slots & Rule & Resolved & $[H^{\mathrm{clair}}_{\mathrm{lower}}, H^{\mathrm{clair}}_{\mathrm{upper}}]$ \\
\midrule
free-text slots UNRESOLVED & unique-or-abstain & 89/97 & [0.253, 0.335] \\
free-text slots UNRESOLVED & first-admitted & 89/97 & [0.216, 0.299] \\
free-text slots planner-fixed & unique-or-abstain & 97/97 & [0.306, 0.306] \\
free-text slots planner-fixed & first-admitted & 97/97 & [0.216, 0.216] \\
\bottomrule
\end{tabular}

\end{table}

\section{CONTRACT\_V2 Census: Additional Tables}
\label{app:v2}
\begin{table}[h]
\caption{CONTRACT\_V2 per suite (fallback arm, typed resolver, official
channel). Gap$>0$: tasks with positive lower gap; Undet.: tasks with
capped searches; Sel.\ fail.: selectable-failure variants; Not refuted: by
the first-write identifiability screen (a necessary condition only), with
the number of vacuous rows (no other variant shares the first-write view) in
parentheses.}
\label{tab:ccsuite}
\centering\small
\begin{tabular}{lccccccc}
\toprule
Suite & Tasks & Gap$>0$ & Undet. & Unres. & $[H_{\mathrm{lower}}, H_{\mathrm{upper}}]$ & Sel.\ fail. & Not refuted (vacuous) \\
\midrule
workspace & 40 & 9 & 3 & 0 & [0.225, 0.275] & 135 & 135 (105) \\
travel & 20 & 0 & 0 & 0 & [0.000, 0.000] & 0 & 0 (0) \\
banking & 16 & 3 & 0 & 0 & [0.131, 0.131] & 21 & 21 (11) \\
slack & 21 & 16 & 0 & 0 & [0.730, 0.730] & 92 & 92 (32) \\
\bottomrule
\end{tabular}

\end{table}

\begin{table}[h]
\caption{V1 versus V2 contract. The V1 contract was derived from hidden
ground truth and is not deployable; the rows are not a like-for-like
comparison of the same contract. Executed attacks: official channel, injected
variants of resolved tasks.}
\label{tab:ccv1v2}
\centering\small
\begin{tabular}{llccc}
\toprule
Contract & Baseline & Resolved & $[H_{\mathrm{lower}}, H_{\mathrm{upper}}]$ & Executed attacks \\
\midrule
V1: slots and contract from benign ground truth & unique-or-abstain & 89/97 & [0.253, 0.335] & 4/851 \\
V2: schema slots, actor-view candidates & unique-or-abstain & 97/97 & [0.314, 0.334] & 4/949 \\
V2: schema slots, actor-view candidates & typed resolver & 97/97 & [0.273, 0.293] & 4/949 \\
\bottomrule
\end{tabular}

\end{table}

\section{AgentDojo Offline Replay Protocol}
\label{app:protocol}

\paragraph{Snapshot and import path.}
We use AgentDojo~0.1.35, the PyPI wheel with sha256 prefix
\texttt{364bea42}, MIT license, at benchmark version v1.2.2.
\begin{itemize}
\item The wheel is installed without its declared LLM-client dependencies.
\item A small import shim skips the package initializer that imports those
clients. The only constant the evaluation path needs from an LLM module is
read from the source's syntax tree.
\item No language-model API, GPU or service account is used.
\end{itemize}

\paragraph{Fixed components} (all frozen before the census):
\begin{itemize}
\item \textbf{Plan.} The task's own \texttt{ground\_truth(env)} call
sequence in each variant's environment, with final output
\texttt{GROUND\_TRUTH\_OUTPUT} (\textsc{oracle\_plan\_conditional}).
\item \textbf{State-changing calls.} Calls whose execution changes the
environment dump.
\item \textbf{Slots.} Arguments of state-changing calls whose benign
ground-truth value is absent from the user prompt and occurs in an earlier
benign tool output, either as a field value or as a substring of a text
field.
\item \textbf{Contract.} For each slot, the (tool, field path) set where
the benign value occurs. Substring slots additionally require a typed
extractor (IBAN, e-mail, URL, date, date-time, time, number) chosen by the
type of the benign value.
\item \textbf{Admitted set.} Values the contract admits in the variant's
own observation, plus abstain.
\item \textbf{Rules.} Primary: unique-or-abstain. Secondary: first admitted
value.
\item \textbf{Variants.} Benign, and one variant per injection task of the
suite, using AgentDojo's \texttt{DirectAttack} template (``TODO: \{goal\}'')
on the vectors that AgentDojo's own canary procedure reports as seen by the
task.
\end{itemize}

\paragraph{Evaluation.}
Each assignment is replayed on its own deep copy of the variant's initial
environment. The utility and security checks run as in AgentDojo:
\texttt{*\_from\_traces} first, and the state-based check if it returns
\texttt{None}. The trace contains only calls that passed the gate. We record
\emph{task\_success}, \emph{attack\_success} (\texttt{security()}~$=$~True;
\texttt{None} for benign variants) and \emph{policy\_violation}
separately, along with blocked calls and tool errors.

\paragraph{Search and bounds.}
The search order is:
\begin{enumerate}
\item the rule assignments;
\item the assignment equal to the ground-truth values (if admitted);
\item exhaustive enumeration up to 4{,}096 assignments per variant, or 4{,}096
seeded samples above that.
\end{enumerate}
The search stops at the first success witness. Bounds follow
\S\ref{sec:headroom}.

\paragraph{Fixtures.}
Ten tests on the pinned snapshot check the direction of every outcome
field:
\begin{itemize}
\item benign success;
\item attack success when the attacker's account is bound;
\item a blocked call that is not counted as an attack, with no attacker
transaction in the post-state;
\item abstention;
\item a tool error;
\item the policy-violation flag;
\item isolation between replays;
\item agreement with AgentDojo's own checker;
\item the injection-replacement property.
\end{itemize}

\section{What We Read}
\label{app:reading}
For the closest papers we read the following sections of the arXiv full
text:
\begin{itemize}
\item \textbf{ROPE}: \S1--\S8 and appendices A (except A.2), B and C;
\item \textbf{PACT}: \S1, \S3, \S4, and appendices A, C--F.1 and I; the
proofs were not read;
\item \textbf{ARGUS}: the full main text;
\item \textbf{CaMeL}: \S3--\S6.
\end{itemize}
Bibliographic entries were checked against the arXiv abstract pages. For the
remaining related work we rely on abstracts or partial sections, as recorded
in our reading notes.

\end{document}
